\DocumentMetadata{
  lang=en-US,
  pdfstandard=ua-2,
  testphase={phase-III,title,table}
}
\documentclass[letterpaper,11pt]{article}

% Assignment metadata. Replace every placeholder before submission.
\newcommand{\TemplateStatus}{template}
% auto shows review annotations until TemplateStatus is final.
\newcommand{\CommentsMode}{auto}
\newcommand{\HomeworkTitle}{Homework 01}
\newcommand{\HomeworkSubtitle}{Template assignment}
\newcommand{\StudentName}{Student Name}
\newcommand{\StudentId}{Student ID}
\newcommand{\CourseCode}{COURSE 101}
\newcommand{\CourseName}{Course Name}
\newcommand{\InstructorName}{Instructor Name}
\newcommand{\SubmissionDate}{Month DD, YYYY}

\usepackage{theme/ricardo-homework}

% Optional reviewer shorthands:
% \DeclareCommenter{\JD}{JD}

\hypersetup{
  pdftitle={\HomeworkTitle: \HomeworkSubtitle},
  pdfauthor={\StudentName},
  pdfsubject={\CourseCode\space coursework},
  pdfkeywords={homework, \CourseCode, \CourseName}
}

\begin{document}

\makehomeworktitle

\begin{problem}[10 points]{A recurrence and its closed form}
  Consider the sequence defined by
  \[
    a_0 = 1,
    \qquad
    a_n = 2a_{n-1} + 1 \quad \text{for } n \geq 1.
  \]
  \begin{subproblems}
    \item Compute the first four terms.
    \item Prove that $a_n = 2^{n+1}-1$ for every $n \geq 0$.
  \end{subproblems}

  \begin{solution}
    The first four terms are $1$, $3$, $7$, and $15$.
    We prove the closed form by induction.
    The base case is $a_0=1=2^1-1$.
    If $a_k=2^{k+1}-1$, then
    \[
      a_{k+1}=2a_k+1
      =2\left(2^{k+1}-1\right)+1
      =2^{k+2}-1.
    \]
    Therefore, the formula holds for all $n\geq 0$.
  \end{solution}
\end{problem}

\begin{problem}[5 points]{Complexity analysis}
  The following algorithm performs a linear search.
  State its worst-case running time and justify your answer.

  \begin{codebox}[python]{linear_search.py}
def locate(values, target):
    for index, value in enumerate(values):
        if value == target:
            return index
    return None
  \end{codebox}

  \begin{solution}
    In the worst case, the algorithm inspects every element exactly once.
    Its running time is therefore $\Theta(n)$ and its auxiliary space usage is $\Theta(1)$.
    The source and explanation are kept together following the general literate-programming principle described by \citet{knuth1984literate}.
  \end{solution}
\end{problem}

\bibliographystyle{plainnat}
\bibliography{references}

\end{document}
